\documentclass[11pt]{article}
\usepackage{amsmath,amssymb,amsthm}
\usepackage[margin=1in]{geometry}
\newtheorem{theorem}{Theorem}
\newtheorem{lemma}{Lemma}
\begin{document}
\title{V109: A Gate-Capacitated Flow Dichotomy for Essential Ternary MUX}
\author{NC0\_k-Avoid Laboratory}
\date{}
\maketitle

A signed essential MUX output has the form
\[
 g=o\oplus \operatorname{MUX}(s\oplus p_s,a\oplus p_a,b\oplus p_b).
\]
Fixing its output gives two exact binary clauses.  Branch $r$ therefore supplies
an implication
\[
 x_s=\alpha_r\Longrightarrow x_d=\beta_r,
\]
where $\alpha_0=p_s$, $\alpha_1=1\oplus p_s$, and the output target chooses the
arrival phase $\beta_r$ on the selected data variable.

\begin{lemma}[Opposite-phase double cycles]
Suppose two directed branch cycles start at the same selector $v$, use disjoint
output gates, and their first source phases are opposite.  Then a word outside
the circuit range can be constructed in time linear in the two cycle lengths.
\end{lemma}

\begin{proof}
On each cycle choose output targets so that every selected implication arrives
at the source phase of the next selected branch, except that the final branch
arrives at the complement of the initial source phase.  If the initial source
phase is $\alpha$, the cycle contains
\[
 x_v=\alpha\Longrightarrow x_v=1\oplus\alpha,
\]
so every satisfying assignment must set $x_v=1\oplus\alpha$.  Opposite initial
phases therefore force opposite Boolean values on $v$.  Gate disjointness makes
the two target assignments compatible; unused MUX clauses only strengthen the
contradiction.
\end{proof}

Fix two distinct outputs with common selector $v$, choose first branches with
opposite source phases, and let their destinations be $d_0,d_1$.  Remove every
output whose selector is $v$ from the return graph.  Split every remaining
output gate $h$ into $h_{\mathrm{in}}\to h_{\mathrm{out}}$ with capacity one;
all variable-to-gate and gate-to-data arcs have capacity two.  A super-source
sends capacity one to each $d_i$, and the sink is $v$.

\begin{lemma}[Return reachability]
If the original branch graph is strongly connected, each $d_i$ reaches $v$
after all selector-$v$ outputs are removed.
\end{lemma}

\begin{proof}
Take a simple directed path from $d_i$ to $v$ in the original branch graph.
Before its final vertex the path never visits $v$, hence it uses no outgoing
branch whose selector is $v$.
\end{proof}

\begin{theorem}[Flow-or-one-gate-bottleneck]
In the gate-capacitated return network the maximum flow has value one or two.
If it has value two, a missing output is constructible in polynomial time.  If
it has value one, there is one explicit output gate lying on every directed
return path from either $d_0$ or $d_1$ to $v$.
\end{theorem}

\begin{proof}
The previous lemma gives value at least one, while the super-source has total
capacity two.  An integral value-two flow decomposes into two source-to-sink
routes disjoint on the unit-capacity gate edges.  Prepending the two prescribed
first branches yields the cycles of the first lemma.

If the value is one, max-flow/min-cut gives a cut of capacity one.  All internal
non-gate edges have capacity two.  A single super-source edge cannot separate
$v$, since the other destination still has an independent return path.  Hence
the unit cut is a gate split edge $h_{\mathrm{in}}\to h_{\mathrm{out}}$.
Deleting $h$ destroys every return path from both destinations to $v$.
\end{proof}

\begin{theorem}[Strong-SCC structural alternative]
For an essential MUX circuit with $m>n$ whose branch graph is strongly connected,
there is a deterministic polynomial-time procedure that outputs either
\begin{enumerate}
\item a missing-output certificate from two gate-disjoint opposite-phase cycles,
or
\item one output gate dominating both return cones for a repeated selector.
\end{enumerate}
\end{theorem}

\begin{proof}
Because every output has exactly one selector and $m>n$, some input is selector
of at least two outputs.  Choose branches of opposite selector source phases and
apply the flow dichotomy.  Scanning the polynomial number of gate pairs and
branch choices remains polynomial.
\end{proof}

\paragraph{Strict separation.}
For each $k\ge2$, take variables $v,A_1,\ldots,A_k,B_1,\ldots,B_k$, two central
outputs $\operatorname{MUX}(v,A_1,B_1)$, and on each lobe outputs
$\operatorname{MUX}(X_i,X_{i+1},v)$ for $i<k$ and
$\operatorname{MUX}(X_k,v,X_1)$.  Then $n=2k+1$, $m=n+1$.  The branch graph is
strongly connected and two gate-disjoint opposite-phase cycles are immediate.
The support family is Hall-minimal.  No ignored-output set exposes the V108
SCC-separated certificate.  Its V102 backdoor is exactly
\[
 \beta=1+2\lceil k/2\rceil=\Theta(n).
\]

\paragraph{Boundary.}
The one-gate bottleneck alternative is not yet converted into a missing-output
algorithm.  Thus V109 does not prove polynomial-time avoidance for all essential
MUX/bijunctive circuits, unrestricted $NC^0_3$-Avoid, a new general circuit lower
bound, or P versus NP.  Novelty and priority are not established.

\end{document}
